\documentclass[11pt]{amsart}

\usepackage[T1]{fontenc}
\usepackage{lmodern}
\usepackage{microtype}
\usepackage{mathtools,amssymb,amsthm}
\usepackage[hidelinks]{hyperref}

\hypersetup{
  pdftitle={The monomial x0 x1^4 x2^5 is not 2-computable under Definition 3.5 as printed}
}

\newtheorem{theorem}{Theorem}
\newtheorem{lemma}[theorem]{Lemma}
\newtheorem{corollary}[theorem]{Corollary}
\theoremstyle{remark}
\newtheorem*{remark}{Remark}
\theoremstyle{definition}
\newtheorem*{definition}{Definition 3.5 of the source}

\DeclareMathOperator{\rk}{rk}
\DeclareMathOperator{\HF}{HF}
\newcommand{\perpp}{^{\perp}}

\title[Not $2$-computable under Definition~3.5 as printed]
{The Monomial $x_0x_1^4x_2^5$ is Not $2$-Computable under Definition~3.5 as Printed}

\date{}

\subjclass[2020]{15A69, 14N07, 13E10}
\keywords{Waring rank, $e$-computability, apolarity, Hilbert function,
artinian Gorenstein algebra, monomial, Strassen's conjecture}

\begin{document}

\begin{abstract}
Carlini, Catalisano, Chiantini, Geramita and Woo introduced $e$-computability as a tool for
Strassen's additivity conjecture, and asked in their Remark~6.2 whether $G_1=x_0x_1^4x_2^5$ is
$2$-computable. Under their Definition~3.5 as printed, the answer is no. The reason is a
one-line upper bound: if $F\in S_d$ is $e$-computable then
$e\cdot\rk(F)\le N(F)-\sum_{j<e}\HF(T/F\perpp,j)$, where $N(F)=\dim_k T/F\perpp$; for $G_1$ this
reads $60\le 56$. A remark records that the same bound forces $e\le a_0$ for every monomial
$x_0^{a_0}\cdots x_n^{a_n}$ with $0<a_0\le\cdots\le a_n$. Nothing here bears on Strassen's
conjecture itself, and nothing here covers readings of $e$-computability that drop the clause
``generated in degree $e$''.
\end{abstract}

\maketitle

\section{The question}

Let $k$ be algebraically closed of characteristic $0$, let $S=k[x_0,\dots,x_n]$ and
$T=k[X_0,\dots,X_n]$, and let $T$ act on $S$ by $X_i\mapsto\partial/\partial x_i$. For
$0\ne F\in S_d$ put $F\perpp=\operatorname{Ann}(F)\subset T$ and $A_F=T/F\perpp$; this is an
artinian Gorenstein algebra of socle degree $d$, so
\[
 \HF(A_F,j)=\HF(A_F,d-j),
 \qquad
 \HF(A_F,j)>0 \iff 0\le j\le d,
\]
and we write $N(F)=\dim_kA_F=\sum_j\HF(A_F,j)$. Finally $\rk(F)$ is the Waring rank of $F$.

Carlini, Catalisano, Chiantini, Geramita and Woo \cite{CCCGW} introduced the following notion.

\begin{definition}
$F\in S_d$ is \emph{$e$-computable} if there exists an ideal $I\subset T$
\emph{generated in degree $e$} such that
\begin{equation}\label{eq:ecomp}
 \rk(F)=\frac1e\sum_i\HF\bigl(T/(F\perpp:I+(t)),\,i\bigr)
\end{equation}
for general $t\in I_e$.
\end{definition}

\noindent
That is their definition restated in the notation fixed above; verbatim, Definition~3.5 of
\cite{CCCGW} reads ``there exists an ideal $I\subset T$ generated in degree $e$ such that for
general $t\in I_e$ we have
$\rk(F)=\bigl(\tfrac1e\bigr)\sum_{i=0}^{\infty}\HF(T/(F\perpp:I+(t)),i)$''. The clause
``generated in degree $e$'' is load-bearing for everything below; see the Scope paragraph. Their
Remark~6.2, on journal page~387, contains the question we answer:

\begin{quote}
If $G_1=x_0x_1^4x_2^5$ then we showed that $G_1$ is $1$-computable and $\rk(G_1)=30$. But we do
not know if $G_1$ is $2$-computable. Thus we cannot use the theorem to find the rank of $F+G_1$,
although Strassen's conjecture says that the rank should be $25+30$. However, if
$G_2=x_0^3x_1^4x_2^5$, by Proposition~4.2, we know that $G_2$ is $2$-computable and
$\rk(G_2)=30$. Hence $\rk(F+G_2)=25+30=55$.
\end{quote}

\noindent
Proposition~4.2 of \cite{CCCGW} makes $x_0^{a_0}\cdots x_n^{a_n}$, with $0<a_0\le\cdots\le a_n$,
$e$-computable for $1\le e\le(a_0+1)/2$; for $G_1$ that gives only $e=1$. Neither that
proposition nor Corollary~3.4 of \cite{CCCGW} bounds $e$ from above, and their Remark~4.3 asks
whether such monomials are ``$e$-computable for $e$'s different from those described in the two
propositions''.

\begin{theorem}\label{thm:main}
$G_1=x_0x_1^4x_2^5$ is not $2$-computable in the sense of Definition~3.5 of \cite{CCCGW} as
quoted above. Consequently the question of Remark~6.2 of \cite{CCCGW} has a negative answer.
\end{theorem}

\section{The obstruction}

Everything follows from one containment. Write $T_{\ge m}=\bigoplus_{i\ge m}T_i$, and read an
empty sum as $0$.

\begin{lemma}\label{lem:main}
Let $0\ne F\in S_d$, let $e\ge1$, let $I\subset T$ be \emph{any} ideal generated in degree $e$,
and let $t\in T$ be \emph{any} homogeneous element. Then $T_{\ge d-e+1}\subseteq F\perpp:I$, and
\begin{equation}\label{eq:cap}
\begin{aligned}
 \sum_i\HF\bigl(T/(F\perpp:I+(t)),\,i\bigr)
 &\;\le\;\sum_{i=0}^{d-e}\HF(A_F,i)\\
 &\;=\;N(F)-\sum_{j=0}^{e-1}\HF(A_F,j)\;<\;N(F).
\end{aligned}
\end{equation}
\end{lemma}

\begin{proof}
\emph{Step 1.} Since $I=(I_e)$, it suffices for $T_{\ge d-e+1}\subseteq F\perpp:I$ to check that
$hg\in F\perpp$ whenever $\deg h\ge d-e+1$ and $g\in I_e$: then $\deg(hg)\ge d+1$, and every form
of degree $>d$ annihilates $F$.

\emph{Step 2.} The containment is strict over $F\perpp$ when $1\le e\le d+1$, because
$\HF(A_F,d-e+1)>0$ forces $T_{d-e+1}\not\subseteq F\perpp$.

\emph{Step 3.} Put $J=F\perpp:I+(t)$; since $F\perpp$, $I$ and $t$ are homogeneous, so is $J$,
and $T/J$ is graded, so its Hilbert function is defined. By Step~1,
$F\perpp+T_{\ge d-e+1}\subseteq F\perpp:I\subseteq J$, and all three quotients are artinian, so
each Hilbert-function sum \emph{is} a $k$-dimension and
\[
\begin{aligned}
 \sum_i\HF(T/J,i)=\dim_kT/J
 &\;\le\;\dim_kT/(F\perpp:I)
 \;\le\;\dim_kT/(F\perpp+T_{\ge d-e+1})\\
 &\;=\;\sum_{i=0}^{d-e}\HF(A_F,i).
\end{aligned}
\]
The equality in \eqref{eq:cap} is Gorenstein symmetry, and the last strict inequality holds
because $\HF(A_F,0)=1$.
\end{proof}

Beyond homogeneity, nothing in the proof uses any property of $t$; in \eqref{eq:ecomp} one has
$t\in I_e$, so $t$ is homogeneous of degree $e$. Hence \eqref{eq:cap} holds under either reading
of ``for general $t\in I_e$'' --- a generic $t$, or merely some $t$ --- and the following
corollary is insensitive to that choice.

\begin{corollary}\label{cor:bound}
If $F$ is $e$-computable then
$e\cdot\rk(F)\le N(F)-\sum_{j=0}^{e-1}\HF(A_F,j)<N(F)$.
\end{corollary}

\begin{proof}[Proof of Theorem~\ref{thm:main}]
Here $d=10$ and $G_1\perpp=(X_0^2,X_1^5,X_2^6)$, the apolar ideal of a monomial: a monomial
$X_0^{c_0}X_1^{c_1}X_2^{c_2}$ kills $x_0x_1^4x_2^5$ exactly when $c_0>1$ or $c_1>4$ or $c_2>5$.
It is a monomial complete intersection, so $A_{G_1}$ has
\[
 \bigl(\HF(A_{G_1},j)\bigr)_{j=0}^{10}=(1,3,5,7,9,10,9,7,5,3,1),
 \qquad
 N(G_1)=2\cdot5\cdot6=60 .
\]
By the monomial Waring-rank theorem of Carlini, Catalisano and Geramita \cite{CCG},
$\rk(G_1)=(4+1)(5+1)=30$, the same value \cite{CCCGW} prints on page~387. If $G_1$ were
$2$-computable, Corollary~\ref{cor:bound} at $e=2$ would give
\[
 60=2\cdot30=e\cdot\rk(G_1)\;\le\;N(G_1)-\HF(A_{G_1},0)-\HF(A_{G_1},1)=60-1-3=56,
\]
which is false, with margin $4$. Hence $G_1$ is not $2$-computable.
\end{proof}

\begin{remark}
The same bound also caps the right-hand side of \eqref{eq:ecomp} for $G_1$ at $e=2$ by $56$, so
the best lower bound for $\rk(G_1)$ obtainable at $e=2$ from Corollary~3.4 of \cite{CCCGW} is
$\lfloor56/2\rfloor=28<30$. Applied to $F=x_0^{a_0}\cdots x_n^{a_n}$ with
$0<a_0\le\cdots\le a_n$, for which $F\perpp=(X_0^{a_0+1},\dots,X_n^{a_n+1})$,
$N(F)=\prod_i(a_i+1)$ and $\rk(F)=N(F)/(a_0+1)$ by \cite{CCG}, Corollary~\ref{cor:bound}
likewise forces $e\le a_0$ whenever $F$ is $e$-computable: for $e\le d+1$ it gives
$e(\rk(F)+1)\le(a_0+1)\rk(F)$, hence $e<a_0+1$, and for $e>d+1$ the left side of \eqref{eq:cap}
is the empty sum $0<e\cdot\rk(F)$. No converse is claimed.
\end{remark}

\section{Verification}

Every number above is a small exact count, and the proof of Theorem~\ref{thm:main} can be redone
by hand from the exponent vector $(1,4,5)$. The companion program \texttt{verify.py} (Python,
standard library only, exact integer and \texttt{Fraction} arithmetic) re-derives it
mechanically: it recomputes $G_1\perpp$ from the apolarity pairing and confirms it equals the
printed ideal on every monomial of degree $\le13$; obtains $\HF(A_{G_1})$ twice, by counting
standard monomials and by ranks of the apolarity matrices; checks the three steps of
Lemma~\ref{lem:main} at $e=2$; and evaluates the cap $56$ as a dimension. It quotes
$\rk(G_1)=30$ from \cite{CCG} rather than reproving it. The transcript also records finite
censuses and control computations that bear on statements no longer made here; those are finite
checks and establish nothing universal. Lemma~\ref{lem:main} itself, which is a statement about
arbitrary $F$, $e$, $I$ and $t$, is proved by hand above and is not established by any finite
computation. All $55$ of its checks pass; the recorded transcript is
\texttt{verify.output.txt}.

\medskip
\noindent\textbf{Scope.} The clause ``generated in degree $e$'' in Definition~3.5 is what makes
Step~1 of Lemma~\ref{lem:main} work. If $I$ is allowed to contain forms of degree $<e$, Step~1
fails and nothing here settles Remark~6.2; restatements of the notion that do not carry that
clause are in circulation, in the authors' withdrawn preprint arXiv:1502.01107, in the survey
\cite{HG}, and in \cite{BCDM}, where it appears instead as a hypothesis of Theorem~2.2. So what
is settled is Theorem~\ref{thm:main} under Definition~3.5 as printed in \cite{CCCGW}. Lemma~\ref{lem:main} and
Corollary~\ref{cor:bound} are characteristic-free, but $\rk(G_1)=30$ is a characteristic-zero
theorem \cite{CCG}, so Theorem~\ref{thm:main} is stated over a field of characteristic $0$. The
bound is necessary, not sufficient, so it settles nothing where it does not fire; in particular
it says nothing about the negative $e=1$ results of \cite{CCCGW}. Nothing here determines
$\rk(F+G_1)$ or bears on Strassen's additivity conjecture: Theorem~\ref{thm:main} closes off the
route of Remark~6.2 rather than completing it, and the value $\rk(F+G_2)=25+30=55$ in the
passage quoted above is the assertion of \cite{CCCGW}, not of this note. Statement numbers are
those of the journal version of \cite{CCCGW}.

\medskip
\noindent\textbf{Attribution and adjacent work.} The shape of the argument is the source's own:
Remark~7.3 of \cite{CCCGW} already bounds the same Hilbert-function sum from above in order to
defeat $1$-computability of $w(x^3+y^3+z^3)$. Corollary~3.4 of \cite{CCCGW} uses the sum in
\eqref{eq:ecomp} to bound $\rk(F)$ from below; Lemma~\ref{lem:main} bounds that sum from above by
an invariant of $F$ alone, uniformly in $I$ and in homogeneous $t$.

\begin{thebibliography}{9}

\bibitem{CCCGW}
E.~Carlini, M.~V. Catalisano, L.~Chiantini, A.~V. Geramita, and Y.~Woo,
\emph{Symmetric tensors: rank, Strassen's conjecture and $e$-computability},
Ann. Sc. Norm. Super. Pisa Cl. Sci. (5) \textbf{XVIII} (2018), no.~1, 363--390.
\href{https://doi.org/10.2422/2036-2145.201602_004}{doi:10.2422/2036-2145.201602\_004};
preprint \href{https://arxiv.org/abs/1506.03176}{arXiv:1506.03176v2}.

\bibitem{CCG}
E.~Carlini, M.~V. Catalisano, and A.~V. Geramita,
\emph{The solution to the Waring problem for monomials and the sum of coprime monomials},
J. Algebra \textbf{370} (2012), 5--14.
\href{https://doi.org/10.1016/j.jalgebra.2012.07.028}{doi:10.1016/j.jalgebra.2012.07.028}.

\bibitem{BCDM}
M.~Bhat, E.~Carlini, S.~Dubey, and S.~K. Masuti,
\emph{Waring decompositions of the product of two quadrics: the small rank cases},
Linear Algebra Appl.,
\href{https://doi.org/10.1016/j.laa.2026.06.027}{doi:10.1016/j.laa.2026.06.027};
\href{https://arxiv.org/abs/2511.23035}{arXiv:2511.23035}.

\bibitem{HG}
A.~Bernardi, E.~Carlini, M.~V. Catalisano, A.~Gimigliano, and A.~Oneto,
\emph{The Hitchhiker guide to: secant varieties and tensor decomposition},
Mathematics \textbf{6} (2018), no.~12, art.~314.
\href{https://doi.org/10.3390/math6120314}{doi:10.3390/math6120314};
\href{https://arxiv.org/abs/1812.10267}{arXiv:1812.10267}.

\end{thebibliography}

\end{document}
